% !TeX program = xelatex
\documentclass[12pt,a4paper]{article}
\usepackage[margin=2cm]{geometry}
\usepackage{fontspec}
\usepackage{amsmath,amssymb}
\usepackage{enumitem}
\setmainfont{Times New Roman}
\setlength{\parindent}{0pt}
\begin{document}
\begin{center}
{\Large\bfseries ĐÁP ÁN ĐỀ THI VÀO LỚP 10 HÀ NỘI NĂM 2007}\\[2pt]
{\bfseries Môn: TOÁN -- PHẦN ĐẠI SỐ}
\end{center}
\textbf{Câu 1.}
\begin{enumerate}[label=\alph*)]
  \item Điều kiện: $x\ge 0$, $x\ne 1$.
  \begin{align*}
  P&=\frac{\sqrt{x}(\sqrt{x}+1)+3(\sqrt{x}-1)-6\sqrt{x}+4}{x-1}\\
   &=\frac{x-2\sqrt{x}+1}{x-1}
    =\frac{\sqrt{x}-1}{\sqrt{x}+1}.
  \end{align*}
  \item
  \[
  P<\frac12
  \iff 2(\sqrt{x}-1)<\sqrt{x}+1
  \iff \sqrt{x}<3
  \iff x<9.
  \]
  Kết hợp điều kiện: $0\le x<9$ và $x\ne 1$.
\end{enumerate}
\textbf{Câu 2.}

Gọi vận tốc lúc đi là $x\,\mathrm{km/h}$, $x>0$. Vận tốc lúc về là $(x+4)\,\mathrm{km/h}$. Theo đề bài:
\[
\frac{24}{x}-\frac{24}{x+4}=\frac12
\iff x^2+4x-192=0
\iff (x-12)(x+16)=0.
\]
Chọn $x=12$. Vậy vận tốc lúc đi là $\boxed{12\,\mathrm{km/h}}$.

\textbf{Câu 3.}
\begin{enumerate}[label=\alph*)]
  \item
  \[
  x^2-3x+2=0\iff(x-1)(x-2)=0.
  \]
  Vậy $x=1$ hoặc $x=2$.
  \item Theo Viète, yêu cầu tương đương với
  \[
  \begin{cases}
  b^2-4c>0,\\
  c=1.
  \end{cases}
  \]
  Suy ra $c=1$ và $b>2$ hoặc $b<-2$.
\end{enumerate}
\textbf{Câu 5.}

Đường thẳng luôn đi qua $A(0;2)$. Khoảng cách từ $O$ tới đường thẳng không vượt quá $OA=2$; dấu bằng xảy ra khi đường thẳng vuông góc với $OA$, tức song song với trục hoành. Vì hệ số góc là $m-1$ nên
\[
m-1=0\iff\boxed{m=1}.
\]
\end{document}
